\documentclass[a4paper,10pt]{article}
\usepackage{tabularx} % extra features for tabular environment
\usepackage{amsmath,amssymb,amsfonts,amsthm}  % improve math presentation
\usepackage{graphicx} % takes care of graphic including machinery
\usepackage{geometry} % decreases margins
\geometry{
  a4paper,
  right=20mm,
  left=20mm,
  top=20mm,
}

%HEADER
%\usepackage{fancyhdr}
%\pagestyle{fancy}
%\lhead{This is my name}
%\rhead{this is page \thepage}
% chktex-file 44
% chktex-file 3
% chktex-file 1
% chktex-file 12
% chktex-file 24
% chktex-file 46
% chktex-file 8
% chktex-file 36

%REFS, CITATI
\usepackage[final]{hyperref} % advs hyper links inside the generated pdf file

\usepackage[backend=biber, style=numeric, sorting=none]{biblatex}
\addbibresource{ref.bib} % <-- Add your .bib file name here

\hypersetup{
	colorlinks=true,       % false: boxed links; true: colored links
	linkcolor=blue,        % color of internal links
	citecolor=blue,        % color of links to bibliography
	filecolor=magenta,     % color of file links
	urlcolor=blue         
}
\renewcommand{\figurename}{Slika}
\renewcommand{\tablename}{Tabela}

%PACKIGI
\usepackage{physics}
\usepackage{mathtools}
\usepackage{bm}
\usepackage{placeins}
\usepackage{tcolorbox}
\usepackage{lipsum}
\usepackage{gensymb}
\usepackage{siunitx}
\usepackage{mdframed}
\usepackage{empheq}
\usepackage{caption}
\usepackage{subcaption}
\usepackage{enumitem}
\usepackage{icomma}
\usepackage{colortbl}
\usepackage{multirow}
\usepackage{dsfont}
\usepackage{fontspec}
\usepackage[version=4]{mhchem}
\usepackage{cuted} 
\usepackage{algorithm}
\usepackage{booktabs}
\usepackage{algpseudocode}
% ---------------------------------------------------------------
\algnewcommand\algorithmicreturnx{\textbf{return}}
\algnewcommand\Ret{\State \algorithmicreturnx{} }
 

%tikz
\usepackage{tikz}
\usepackage[compat=1.1.8]{tikz-feynman}
%jeziki
\usepackage[slovene]{babel}

%FONT SETUP
\setmainfont{Libertinus Serif}

\setsansfont{Schnebel Sans Pro}[
  BoldFont = {Schnebel Sans Pro Bold},
  UprightFont = {Calibri}]

%CAPTION SETUP
\usepackage[labelfont=bf, font={sf, footnotesize}, labelformat=simple]{caption}
\DeclareCaptionFormat{nature}{\textbf{#1#2}~\footnotesize#3}
\DeclareCaptionLabelSeparator{pipe}{\,$\bm{|}$}

\captionsetup
{
    format=nature,
    labelsep=pipe
}

%SECTION/SUBSECTION SETUP
\usepackage[sc]{titlesec}
% \titleformat*{\section}{\Large\scshape\bfseries\centering}
% \titleformat*{\subsection}{\large\scshape\bfseries\centering}
% \titleformat*{\subsubsection}{\subsubsectionfont\centering}
\titleformat*{\section}{\Large\bfseries\sffamily\centering}
\titleformat*{\subsection}{\large\bfseries\sffamily\centering}
\titleformat*{\subsubsection}{\itshape\subsubsectionfont\centering}


%MISC
\usepackage{xprintlen}
\DeclareMathOperator\arctanh{arctanh}

\definecolor{titlegray}{gray}{0.9}
\definecolor{lightgray}{gray}{0.92} 

\renewcommand{\vec}[1]{\mathbf{#1}}
\AtBeginDocument{\RenewCommandCopy\qty\SI}
%________________________________________________________________________________
\begin{document}
% \begin{titlepage}
% \scshape

% \centering
% \includegraphics[width=12cm, keepaspectratio]{UL_FMF-logoENG-HOR-CMYK_color.pdf} \par \vspace{0.5cm}



% \raisebox{-\baselineskip}{\rule{\textwidth}{1pt}}
% \rule{\textwidth}{1pt}
% \vspace{0.2cm}

%     {\fontsize{24pt}{26pt}\selectfont\bfseries Teorija trdne snovi} \par \vspace{0.1cm}
% {\fontsize{14pt}{20pt}\selectfont Peierlsova nestabilnost}

% \raisebox{-\baselineskip}{\rule{\textwidth}{1pt}}
% \rule{\textwidth}{1pt}
% \vspace{0.2cm}
% \vspace{1cm}


% {\fontsize{12pt}{18pt}\selectfont
% Korl Železnikar\\ \vspace{0.2cm}
% Mentorja: doc. dr. Denis Golež}

% \vspace{2cm}




% \vfill
% {\fontsize{12pt}{18pt}\selectfont \today}

% \end{titlepage}

\section{Osnove teorije}
\subsection{BCS teorija pri končni temperaturi}
BCS model, kjer je $\xi_{\vec{k}, \sigma} = \epsilon_{\vec{k}, \sigma} - E_F$ in $k := (\vec{k}, \uparrow)$, $-k := (\vec{-k}, \downarrow)$:
\begin{align}
  H = \sum_k \xi_k (c_k^\dagger c_k + c_{-k}^\dagger c_{-k}) - V \sum_{k,k'} c_k^\dagger c_{-k}^\dagger c_{-k'} c_{k'}
\end{align}
Naredimo približek povprečnega polja, definiramo $b_k = c_{k} c_{-k}$ in $b_k = \langle b_k\rangle + \delta b_k = F_k  + \delta b_k$, sledi (do $\mathcal{O}(\delta b_k^2)$)
\begin{align}
  b_k^\dagger b_{k'} \simeq  F_k^* F_{k'} + F_k^* \delta b_{k'} + F_{k'} \delta b_k^\dagger = F_k^* b_{k'} + F_{k'} b_k^\dagger - F_k^* F_{k'} =  \langle b_k^\dagger \rangle b_{k'} + b_k^\dagger \langle b_{k'} \rangle - \langle b_k^\dagger \rangle \langle b_{k'} \rangle 
\end{align}
Iz česar sledi: 
\begin{align}
  H = \sum_k \begin{pmatrix} c_k^\dagger & c_{-k} \end{pmatrix} \begin{pmatrix} \xi_k & -\Delta_k \\ -\Delta_k^* & -\xi_k \end{pmatrix} \begin{pmatrix} c_k \\ c_{-k}^\dagger \end{pmatrix} + \frac{\vert{}\Delta\vert{}^2}{V} + \sum_k \xi_k
\end{align}
Ob definiciji $\Delta_k = \sum_{k'} V(k-k') \langle c_{k'}c_{-k'}\rangle$, vzamemo še lokalno interakcijo, zato je $V(q)\simeq V$ in $\Delta_k \simeq \Delta$. Naredimo Bogoljubov rotacijo, privzamemo $u_k, v_k \in \mathbb{R}$ in simetrijo okoli $E_F$: $u_k = u_{-k}$ ter $v_k = v_{-k}$. Iz  CCR sledi $u_k^2+v_k^2=1$. 
Kvazidelci so definirani z rotacijo 
\begin{align}
  \begin{pmatrix} c_k \\ c_{-k}^\dagger \end{pmatrix} = \begin{pmatrix} u_k & v_k \\ -v_k & u_k \end{pmatrix} \begin{pmatrix} \alpha_k \\ \alpha_{-k}^\dagger \end{pmatrix}
\end{align}
nov Hamiltonjan pa postane
\begin{align}
  H = \sum_k \begin{pmatrix} \alpha_k^\dagger & \alpha_{-k} \end{pmatrix}   \begin{pmatrix} \xi_k u_k^2 + u_k v_k (\Delta_k + \Delta_k^*) - \xi_k v_k^2 & -2\xi_k u_k v_k + \Delta_k u_k^2 - \Delta_k^* v_k^2 \\ -2\xi_k u_k v_k - \Delta_k v_k^2 + \Delta_k^* u_k^2 & \xi_k v_k^2 - \xi_k u_k^2 + u_k v_k (\Delta_k + \Delta_k^*) \end{pmatrix} \begin{pmatrix} \alpha_k \\ \alpha_{-k}^\dagger \end{pmatrix}\,,
\end{align} 
kjer zahtevamo, da so izvendiagonalni členi enaki 0. Iz CCR zahteve lahko vzamemo parameterizacijo (tu sem vzel kompleksno nastavimo, da je $\Delta = \Delta^* e^{i2\phi_k}$):
\begin{align}
v_k &= \sin\left(\frac{\theta_k}{2}\right) e^{i\phi_k} \\
u_k &= \cos\left(\frac{\theta_k}{2}\right) \\
u_k v_k &= \cos\left(\frac{\theta_k}{2}\right) \sin\left(\frac{\theta_k}{2}\right) = \frac{1}{2} \sin\theta_k\\
u_k^2 - v_k^2 &= \cos^2\left(\frac{\theta_k}{2}\right) - \sin^2\left(\frac{\theta_k}{2}\right) = \cos\theta_k 
\end{align}
skupaj dobimo (od tu naprej $\Delta := |\Delta|$): 
\begin{align}
  \xi_k \sin(\theta_k) =  \Delta \cos(\theta_k)
\end{align}
kar fiksira 
\begin{align}
  \tan(\theta_k) = \frac{|\Delta_k|}{\xi_k}\,.
\end{align}
Tako iz trikotnika preberemo
\begin{align}
  \sin\theta_k = \frac{\vert{}\Delta_k\vert{}}{E_k} = \frac{\vert{}\Delta_k\vert{}}{\sqrt{\xi_k^2 + \vert{}\Delta_k\vert{}^2}} \\
  \cos\theta_k = \frac{\xi_k}{E_k} = \frac{\xi_k}{\sqrt{\xi_k^2 + \vert{}\Delta_k\vert{}^2}} \\
\end{align}
kjer je $E_k = \pm\sqrt{\xi_k^2 + \vert{}\Delta_k\vert{}^2}$ energija kvazidelcev:
\begin{align}
  H = \sum_k E_k \alpha_k^\dagger \alpha_k - E_k \alpha_{-k}^\dagger \alpha_{-k} + \xi_k +  \frac{|\Delta|^2}{V}\,.
\end{align} 
V termičnem stanju so diagonalni delci porazdeljeni po 
\begin{align}
  \langle \alpha_k^\dagger \alpha_k\rangle = \frac{1}{1+ \exp(E_k / T)} = \frac{1}{2}-\frac{1}{2} \tanh\left(\frac{E_k}{2T}\right)\,.
\end{align}
Uporabni sta zvezi $f(-E/T) = 1-f(E/T)$ in $1-2f(E/T) = \tanh(E/(2T))$. V definiciji ureditvenega parametra je 
\begin{align}
  \langle c_k c_{-k} \rangle = \left\langle \left( u_k \alpha_k - v_k \alpha_{-k}^\dagger \right) \left( u_k \alpha_{-k} + v_k \alpha_k^\dagger \right) \right\rangle= u_k v_k \left( 1 - \langle \alpha_{-k}^\dagger \alpha_{-k} \rangle - \langle \alpha_k^\dagger \alpha_k \rangle \right) = \frac{1}{2} \sin\theta_k \tanh\left(\frac{\beta E_k}{2}\right)
\end{align}
iz česar sledi enačba za režo pri $T>0$ (seštevamo po $k$ okoli $E_F$):
\begin{align}
  \Delta = \frac{V}{2} \sum^{'}_k \sin\theta_k \tanh\left(\frac{E_k}{2T}\right) = \frac{V \Delta}{2} \sum_k \frac{\tanh\left(\frac{1}{2T} \sqrt{\Delta^2 + \xi_k^2}\right)}{\sqrt{\Delta^2 + \xi_k^2}}
\end{align}
\subsection{Minimalni model ekscitonskega izolatorja pri končni temperaturi}
Pri izolatorju za minimalni model vzamemo dve orbitali ($a, b$) model tesne vezi v vsaki posebej in lokalno Coulombovo interakcijo na mestu $j$ med elektronskima gostotama v posameznih pasovih 
\begin{align}
  H = H_0 + H_V &= \sum_{i,j} \left( \epsilon_a^{(0)} \delta_{ij} + t_{ij}^{(a)} \right) c_{a,i}^\dagger c_{a,j} + \sum_{i,j} \left( \epsilon_b^{(0)} \delta_{ij} + t_{ij}^{(b)} \right) c_{b,i}^\dagger c_{b,j} + V \sum_j c_{a,j}^\dagger c_{a,j} c_{b,j}^\dagger c_{b,j} \\
  &= \sum_k \epsilon_{a,k} c_{a,k}^\dagger c_{a,k} + \sum_k \epsilon_{b,k} c_{b,k}^\dagger c_{b,k} + V \sum_j c_{a,j}^\dagger c_{a,j} c_{b,j}^\dagger c_{b,j}
\end{align}
najprej diagonaliziramo vsako orbitalo posebej, kar da gole disperzije (vezavni energiji nastavimo na $\pm G/2$) (delamo v dveh dimenzijah in mrežna konstanta $a_0=1$): 
\begin{align} \label{free el tb}
  \epsilon_{a(b),k} = \pm \frac{G}{2} - 2 t_{a(b)}(\cos(k_x) + \cos(k_y))\,,
\end{align}
v tem primeru je najmanjša razdalja med pasovoma $G_0 = G - 4 (t_a -t_b)$ oziroma $G_0=G - 8t_a$, če je $t_b = -t_a$. Za interakcijo spet naredimo približek povprečnega polja: 
\begin{align}
  c_{a,j}^\dagger c_{a,j} c_{b,j}^\dagger c_{b,j} \approx \langle c_{a,j}^\dagger c_{a,j} \rangle c_{b,j}^\dagger c_{b,j} + \langle c_{b,j}^\dagger c_{b,j} \rangle c_{a,j}^\dagger c_{a,j} - \langle c_{a,j}^\dagger c_{b,j} \rangle c_{b,j}^\dagger c_{a,j} - \langle c_{b,j}^\dagger c_{a,j} \rangle c_{a,j}^\dagger c_{b,j} + \text{konst.}
\end{align}
od tod sledi
\begin{align}
  H_V^{\text{MF}} = V n_b \sum_j c_{a,j}^\dagger c_{a,j} + V n_a \sum_j c_{b,j}^\dagger c_{b,j} - V \phi \sum_j c_{b,j}^\dagger c_{a,j} - V \phi^* \sum_j c_{a,j}^\dagger c_{b,j} + N V \vert{}\phi\vert{}^2
\end{align}
kjer smo definirali: 
\begin{align}
  n_a = \langle c_{a,j}^\dagger c_{a,j} \rangle = \frac{1}{N} \sum_k \langle c_{a,k}^\dagger c_{a,k} \rangle
\end{align}
\begin{align}
  n_b = \langle c_{b,j}^\dagger c_{b,j} \rangle = \frac{1}{N} \sum_k \langle c_{b,k}^\dagger c_{b,k} \rangle
\end{align}
\begin{align}
  \phi = \langle c_{a,j}^\dagger c_{b,j} \rangle = \frac{1}{N} \sum_k \langle c_{a,k}^\dagger c_{b,k} \rangle \\
  \Delta = -\frac{V}{N} \sum_k \langle c_{a,k}^\dagger c_{b,k} \rangle, \quad \Delta^* = -\frac{V}{N} \sum_k \langle c_{b,k}^\dagger c_{a,k} \rangle
\end{align}
\begin{align}
  \Delta_k = \frac{1}{N} \sum_{k'} V(k - k') \langle a_{k'}^\dagger b_{k'} \rangle, \quad V(q) = \int e^{-i \vec{q} \cdot \vec{r}} V(\vec{r}) \, \dd{\vec{r}}
\end{align}
Uporabimo FT še na celotnem Hamiltonjanu
\begin{align}
  c_{a,j}^\dagger = \frac{1}{\sqrt{N}} \sum_k e^{-i \vec{k} \cdot \vec{r}_j} c_{a,k}^\dagger
\end{align}
\begin{align}
  H^{\text{MF}} = \sum_k \left( [\epsilon_{a,k} + V n_b] c_{a,k}^\dagger c_{a,k} + [\epsilon_{b,k} + V n_a] c_{b,k}^\dagger c_{b,k} + \Delta^* c_{a,k}^\dagger c_{b,k} + \Delta c_{b,k}^\dagger c_{a,k} \right) + \left( \frac{\vert{}\Delta\vert{}^2}{V} N - V N n_a n_b \right)
\end{align}
oziroma: 
\begin{align}
  H^{\text{MF}} = \sum_k \begin{pmatrix} c_{a,k}^\dagger & c_{b,k}^\dagger \end{pmatrix} \begin{pmatrix} {\epsilon}^H_{a,k} & \Delta^* \\ \Delta & {\epsilon}^H_{b,k} \end{pmatrix} \begin{pmatrix} c_{a,k} \\ c_{b,k} \end{pmatrix} + \text{konst.}\,,
\end{align}
kjer definiramo Hartree energiji ${\epsilon}^H_{a(b),k} = \epsilon_{a(b),k} + V n_{b(a)}$. Transformacija v diagonalno  / Bogoljubov bazo je podobna kot pri BCS, le da ne meša kreatorjev in anihilatorjev / ohranja število delcev:
\begin{align}
  \begin{pmatrix} c_{a,k} \\ c_{b,k} \end{pmatrix} = \begin{pmatrix} u_k & v_k \\ -v_k & u_k \end{pmatrix} \begin{pmatrix} \alpha_k \\ \beta_k \end{pmatrix}
\end{align}
od tod po enakem postopku kot prej dobimo energiji kvazidelcev
\begin{align}
E_{\alpha(\beta)}(k) = \frac{{\epsilon}^H_{a,k} + {\epsilon}^H_{b,k}}{2} \pm \sqrt{\left(\frac{{\epsilon}^H_{a,k} - {\epsilon}^H_{b,k}}{2}\right)^2 + \vert{}\Delta\vert{}^2} := \eta_k \pm \sqrt{\xi_k^2 +\Delta^2}
\end{align}
Za bookkeeping: 
\begin{align}
  \left( \frac{\epsilon_a^H - \epsilon_b^H}{2} \right)& \sin \theta_k + \Delta \cos \theta_k = 0 \implies \tan \theta_k = \frac{-2 \Delta}{\epsilon_a^H - \epsilon_b^H} = \frac{-\Delta}{\xi_k}\\
  u_k v_k &= -\frac{\sin \theta_k}{2}\,.
\end{align}
Za samousklajeno enačbo zdaj potrebujemo
\begin{align}
  \langle c_{a,k}^\dagger c_{b,k} \rangle = \langle (u_k \alpha_k^\dagger + v_k \beta_k^\dagger)(-v_k \alpha_k + u_k \beta_k) \rangle = u_k v_k \left( \langle \beta_k^\dagger \beta_k \rangle - \langle \alpha_k^\dagger \alpha_k \rangle \right) = -\frac{1}{2} \sin\theta_k (f(E_{\beta}) - f(E_{\alpha}))
\end{align}
od tod
\begin{align}
  \Delta = -\frac{V}{N}\sum_k \langle c_{a,k}^\dagger c_{b,k} \rangle = \frac{V}{N} \sum_k \frac{ \Delta}{2\sqrt{\xi_k^2 + \vert{}\Delta\vert{}^2}} (f(E_{\beta}) - f(E_{\alpha}))
\end{align}
Opomba, lahko definiramo tudi $n = n_a + n_b = \text{konst.} = 1$ in $n_a-n_b =: m$ od tod sledi

\begin{align}
  H^{\text{MF}} = \sum_k \left( \left[\epsilon_{a,k} + \frac{V}{2}n - \frac{V}{2}m\right] c_{a,k}^\dagger c_{a,k} + \left[\epsilon_{b,k} + \frac{V}{2}n -\frac{V}{2}m\right] c_{b,k}^\dagger c_{b,k} + \Delta^* c_{a,k}^\dagger c_{b,k} + \Delta c_{b,k}^\dagger c_{a,k} \right)
\end{align}
od koder je v 2D Hartree gap enak: 
\begin{align}
  G^H = G-4(t_a-t_b) - V m\,. 
\end{align}
Note: v praksi je $m<0$
V dveh dimenzijah se Hartree pasova srečata na 
\begin{align}
  \mu_B = \frac{V}{2}n - \frac{t_a + t_b}{t_a-t_b}\frac{G-Vm}{2}\,.
\end{align}
Zamik je enak $V/2$, ko sta osnovna pasova simetrična in je zasedenost polovična. V primeru simetričnih pasov sta disperziji kvazidelcev
\begin{align}
  E_{\alpha(\beta)}(k) - V/2 = \pm \sqrt{\left(\frac{{\epsilon}_{a,k} - {\epsilon}_{b,k} - Vm}{2}\right)^2 + \vert{}\Delta\vert{}^2}\,.
\end{align}
Velja še
\begin{align}
  n_a &= \frac{1}{N} \sum_k \left( u_k^2 f(E_\alpha) + v_k^2 f(E_{\beta}) \right) \\
  n_b &= \frac{1}{N} \sum_k \left( v_k^2 f(E_\alpha) + u_k^2 f(E_{\beta}) \right)
\end{align}

\begin{align}
  u_k^2 &= \frac{1}{2}\left(1 + \frac{\xi_k}{\sqrt{\Delta^2 + \xi_k^2}}\right) \\
  v_k^2 &= \frac{1}{2}\left(1 - \frac{\xi_k}{\sqrt{\Delta^2 + \xi_k^2}}\right)
\end{align}
\subsection{Numerika/psevdokoda}

\subsection{Odprti sistemi in modeliranje kopeli}
Born-Markov master enačbe in Redfield: ločimo kopel in sistem, 
\begin{align}
  \rho_{S}(t) = \text{Tr}_{B}(\rho(t))
\end{align}
Hamiltonjan razcepimo na $H = H_0 + H_I$, kjer je $H_0 = H_S \otimes \mathbb{I}_B + \mathbb{I}_S \otimes H_B$ razdenjen na podsistema, za interakcijski del pa privzamemo obliko $H_I = \sum_\alpha A_\alpha \otimes B_\alpha$. V interakcijski sliki je operator za evolucijo
\begin{align}
  U_{0}(t, t_{0}) = e^{-i H_{0}(t - t_{0})}
\end{align}
\begin{align}
  \rho_I(t) &= U_0^\dagger(t, 0) \rho(t) U_0(t, 0) = e^{i H_0 t} \rho(t) e^{-i H_0 t}\\
  H_I(t) &= U_0^\dagger(t, 0) H_I U_0(t, 0) = e^{i H_0 t} H_I e^{-i H_0 t}
\end{align}
Von Neumannova enačba je 
\begin{align}
  \frac{d}{dt}\rho_I(t) &= -i [H_I(t), \rho_I(t)] \\
  \rho_I(t) &= \rho_I(0) - i \int_{0}^{t} dt' \, [H_I(t'), \rho_I(t')]
\end{align}
vstavimo nazaj v komutator in se znebimo prostosnih stopenj kopeli:
\begin{align}
  \frac{d}{dt}\rho_{S}(t) = -\int_{0}^{t} ds \, \text{Tr}_{B}\left( [H_{I}(t), [H_{I}(t-s), \rho_{S}(t) \otimes \rho_{B}]] \right)
\end{align}
V interakcijskem členu so operatorji $A_\alpha$ napisani kot lastni operatorji $H_S$ z energijami prehodov $\omega = E_b - E_a$
\begin{align}
  A_\alpha(\omega) = \sum_{E_b - E_a = \omega} \vert{}a\rangle \langle a\vert{} A_\alpha \vert{}b\rangle \langle b\vert{}
\end{align}
To zadošča relacijam $[H_S, A_\alpha(\omega)] = -\omega A_\alpha(\omega), \quad [H_S, A_\alpha^\dagger(\omega)] = \omega A_\alpha^\dagger(\omega)$ in \\ $e^{i H_S t} A_\alpha(\omega) e^{-i H_S t} = e^{-i \omega t} A_\alpha(\omega), \quad e^{i H_S t} A_\alpha^\dagger(\omega) e^{-i H_S t} = e^{i \omega t} A_\alpha^\dagger(\omega)$. V interakcijski sliki potem velja
\begin{align}
  H_I(t) = \sum_{\alpha, \omega} e^{-i \omega t} A_\alpha(\omega) \otimes B_\alpha(t) = e^{i H_B t} B_\alpha e^{-i H_B t} 
\end{align}
evolucija sledi kot
\begin{align}
  \dot{\rho}_{S} = \sum_{\omega, \omega'} \sum_{\alpha, \beta} e^{i(\omega' - \omega)t} \Gamma_{\alpha\beta}(\omega) \left( A_{\alpha}(\omega) \rho_{S}(t) A_{\beta}^{\dagger}(\omega') - A_{\beta}^{\dagger}(\omega') A_{\alpha}(\omega) \rho_{S}(t) \right) + \text{h.c.}
\end{align}
kjer so prehodne matrike med stanji
\begin{align}
  \Gamma_{\alpha\beta}(\omega) = \int_{0}^{\infty} ds \, e^{i\omega s} \, \text{Tr}_{B}\left( B_{\alpha}^{\dagger}(t) B_{\beta}(t-s) \rho_{B} \right) = \frac{1}{2}\gamma_{\alpha\beta}(\omega) + i S_{\alpha\beta}(\omega)
\end{align}
Evolucija pa ima strukturo
\begin{align}
  \frac{d\rho_{S}}{dt} = -i [H_{LS}, \rho_{S}(t)] + \mathcal{D}(\rho_{S}(t))
\end{align}
z disipatorjem
\begin{align}
  \mathcal{D}(\rho_{S}) = \sum_{\omega} \sum_{\alpha,\beta} \gamma_{\alpha\beta}(\omega) \left( A_{\beta}(\omega) \rho_{S} A_{\alpha}^{\dagger}(\omega) - \frac{1}{2} \{ A_{\alpha}^{\dagger}(\omega) A_{\beta}(\omega), \rho_{S} \} \right)
\end{align}
oz. v splošnem
\begin{align}
  \mathcal{D}(\rho_{S}) = \sum_{k} \gamma_{k} \left( L_{k} \rho_{S} L_{k}^{\dagger} - \frac{1}{2} \{ L_{k}^{\dagger} L_{k}, \rho_{S} \} \right)
\end{align}
V stacionarnem stanju velja
\begin{align}
  \mathcal{L}\rho_{\text{ness}} = 0\,.
\end{align}
Notes za vektorizacijo in superoperatorje: 
\begin{align}
\mathcal{L} u_{\mu}^R = \lambda_{\mu} u_{\mu}^R
\end{align}  
\begin{align}
  \text{vec}(ABC) = (A \otimes C^T) \text{vec}(B)
\end{align}
\begin{align}
  \text{Tr}(A^{\dagger} B) = \text{vec}(A)^{\dagger} \text{vec}(B)
\end{align}
\begin{align}
  \frac{d}{dt}\vert{}\rho\rangle = \left( -i H \otimes \mathbb{I} + i \mathbb{I} \otimes H^T + \sum_{k} L_{k} \otimes L_{k}^* - \frac{1}{2} \sum_{k} (L_{k}^{\dagger} L_{k} \otimes \mathbb{I} + \mathbb{I} \otimes L_{k}^T L_{k}^*) \right) \vert{}\rho\rangle
\end{align}
\begin{align}
  \langle \hat{O} \rangle(t) = \langle \mathbb{I} \vert{} \hat{O} \otimes \mathbb{I} \vert{} \rho(t) \rangle
\end{align}
Šibke sklopitve in perturbacijski pristopi
\begin{align}
  \mathcal{L} = -i [H, \cdot] + \epsilon \mathcal{D}[\cdot], \quad  \quad \rho = \rho_{0} + \epsilon \rho_{1} + \epsilon^2 \rho_{2} + \dots
\end{align}
napišemo gostotno matriko sistema v lastni bazi, siagonalno, s klasičnimi populacijami
\begin{align}
  \rho_{0} = \sum_{n} p_{n} \vert{}n\rangle \langle n\vert{}, \quad H\vert{}n\rangle = E_{n} \vert{}n\rangle
\end{align}
kjer velja v ničtem redu (by construction)
\begin{align}
  \mathcal{L}\rho_0 = -i [H, \rho_0] = 0
\end{align}
V prvem redu so enačbe za NESS
\begin{align}
  -i [H, \rho_1] + \mathcal{D}(\rho_0) = 0
\end{align}
kjer je $\mathcal{D}(\rho_0) = \sum_\mu \left( L_\mu \rho_0 L_\mu^\dagger - \frac{1}{2} \{L_\mu^\dagger L_\mu, \rho_0\} \right)$ in od tod enačba za prehode med populacijami
\begin{align}
  \langle n \vert{} \mathcal{D}(\rho_0) \vert{} n \rangle = \sum_\mu \sum_m p_m \vert{}\langle n \vert{} L_\mu \vert{} m \rangle\vert{}^2 - \sum_\mu \sum_m p_n \vert{}\langle m \vert{} L_\mu \vert{} n \rangle\vert{}^2 
\end{align}
kot klasična populacijska dinamika z $R_{nm} = \sum_\mu \vert{}\langle n \vert{} L_\mu \vert{} m \rangle\vert{}^2$
\begin{align}
  \dot{p}_n = \epsilon \sum_m \left( R_{nm} p_m - R_{mn} p_n \right)
\end{align}
Izvendiagonali elementi se razvijejo hitro na $\mathcal{O}(1)$, medtem ko se populacije razvijajo počasi na GGE manifoldu z $\mathcal{O}(\epsilon^{-1})$. Za ohranitev normalizacije mora veljati $\sum_n \dot{p}_n = 0$.

Integrabilni sistemi imajo makrsoskopsko število ohranjenih nabojev $[H_0, C_i]=0$, pod unitarno časovno evolucijo se razvijejo v netermično stanje
\begin{align}
  \rho_{GGE} = \frac{1}{Z} e^{-\sum_{i} \lambda_{i} C_{i}}
\end{align}
\begin{align}
  \langle \psi_0 \vert{} C_i \vert{} \psi_0 \rangle = \text{Tr}(C_i \rho_{\text{GGE}})
\end{align}
fiksira $\lambda_i$. UNitarna dinamika prijcira na GGE manifold, šibka sklopitev premika po GGE manifoldu: 
\begin{align}
  \rho_0(t) = \rho_{\boldsymbol{\lambda}(t)} = \frac{1}{Z(t)} \exp\left( -\sum_i \lambda_i(t) C_i \right)
\end{align} 
EOM za ohranjene naboje je  
\begin{align}
  \langle \dot{C}_{i} \rangle_{\lambda(\infty)} = \text{Tr}[C_{i} \mathcal{D}\rho] \simeq \epsilon \, \text{Tr}[C_{i} \mathcal{D} \rho_{0}] = 0\,\quad \forall \lambda_i 
\end{align} 
Definiramo projektor na GGE manifold:
\begin{align}
  P X = -\sum_{ij} \frac{\partial \rho}{\partial \lambda_{i}} (\chi^{-1})_{ij} \text{Tr}[C_{j} X]
\end{align}   
ki da dinamiko vzdolž manifolda: 
\begin{align}
  \frac{\partial \rho_{\boldsymbol{\lambda}}}{\partial t} = \sum_i \frac{\partial \rho_{\boldsymbol{\lambda}}}{\partial \lambda_i} \dot{\lambda}_i = \epsilon P \mathcal{D}(\rho_{\boldsymbol{\lambda}}) = -\epsilon \sum_{ij} \frac{\partial \rho_{\boldsymbol{\lambda}}}{\partial \lambda_i} (\chi^{-1})_{ij} \text{Tr}\left[ C_j \mathcal{D}(\rho_{\boldsymbol{\lambda}}) \right]
\end{align}
od koder sledi evolucija za koeficiente/potenciale $\lambda_i$:
\begin{align}
  \dot{\lambda}_{i} = -\epsilon \sum_{j} (\chi^{-1})_{ij} \text{Tr}[C_{j} \mathcal{D} \rho_{\lambda(t)}]
\end{align}
s susceptibilnostmi
$\chi_{ij} = \langle C_{i} C_{j} \rangle - \langle C_{i} \rangle \langle C_{j} \rangle = -\text{Tr}\left[ C_{i} \frac{\partial \rho}{\partial \lambda_{j}} \right]$. 
\section{Ekscitonski izolatorji in sklopitev s kopelmi}

\begin{align} 
  \text{Tr}\left[ f\frac{1}{2} \{ \rho, A \} \right] &= \text{Tr}\left[ \frac{1}{2} f \rho A + \frac{1}{2} f A \rho \right] = \frac{1}{2} \text{Tr}\left[ A f \rho + f A \rho \right] \\ &= \frac{1}{2} \text{Tr}\left[ (A f + f A) \rho \right] = \frac{1}{2} \langle A f + f A \rangle 
\end{align}

% \begin{align} 
%   \dot{n}_k = \langle \dot{n}_k \rangle &= \text{Tr}[n_k(t) \rho(t)] = -i \text{Tr}(n_k(t) [H(t), \rho(t)]) + \gamma_0 \text{Tr}[n_k(t) \mathcal{D} \rho] = 0
%  \end{align}

% \begin{align}
% \cdot \quad 2. \quad L_{qp}^{\lambda \lambda'} &= \sqrt{\gamma_{qp}^{\lambda \lambda'}} \, d_{\lambda q}^\dagger d_{\lambda' p} \quad \text{kjer } (\lambda, \lambda') \in \{(\alpha\alpha), (\beta\beta), (\alpha\beta), (\beta\alpha)\} \\
% \Rightarrow \quad \dot{n}_{\alpha k} &= \text{Tr}\left[ n_{\alpha k} \mathcal{D}(\rho) \right] = \sum_{\substack{q,p \\ \lambda, \lambda'}} \text{Tr}\left[ n_{\alpha k} \left( L_{qp}^{\lambda \lambda'} \rho(t) L_{qp}^{\lambda \lambda'\dagger} - \frac{1}{2} \left\{ L_{qp}^{\lambda \lambda'\dagger} L_{qp}^{\lambda \lambda'}, \rho \right\} \right) \right] \\
% &= \sum_{\substack{q,p \\ \lambda, \lambda'}} \langle [L^\dagger, [n, L]] \rangle = ? \\
% &\quad [n_{\alpha k}, d_{\lambda q}^\dagger d_{\lambda' p}] = d_{\lambda q}^\dagger [n_{\alpha k}, d_{\lambda' p}] + [n_{\alpha k}, d_{\lambda q}^\dagger] d_{\lambda' p} \\
% &\quad \text{kjer } d_{\lambda q}^\dagger \{d_{\alpha k}^\dagger, d_{\lambda' p}\} - \{d_{\alpha k}^\dagger, d_{\lambda q}^\dagger\} d_{\lambda' p} = -\delta_{\alpha\lambda'} \delta_{kp} \, d_{\lambda' p} \\
% &\quad = -d_{\lambda q}^\dagger d_{\lambda' p} \, \delta_{\alpha\lambda'} \delta_{kp} + d_{\lambda q}^\dagger d_{\lambda' p} \, \delta_{\alpha\lambda} \delta_{kq} \\
% &= \sum_{\substack{q,p \\ \lambda, \lambda' \\ T_1, T_2}} \left[ \langle d_{\lambda' p}^\dagger d_{\lambda q} (-d_{\lambda q}^\dagger d_{\lambda' p}) \rangle \gamma_{qp}^{\lambda \lambda'} \, \delta_{\alpha\lambda'} \delta_{kp} + \langle d_{\lambda' p}^\dagger d_{\lambda q} d_{\lambda q}^\dagger d_{\lambda' p} \rangle \delta_{\alpha\lambda} \delta_{kq} \, \gamma_{qp}^{\lambda \lambda'} \right] \\
% &= \sum_{p, \lambda'} \Gamma_{kp}^{\alpha \lambda'} n_{\lambda' p} (1 - n_{\alpha k}) - \sum_{q, \lambda} \Gamma_{qk}^{\lambda \alpha} n_{\alpha k} (1 - n_{\lambda q}) \\
% \dot{n}_{\alpha k} &= \sum_{\substack{p, \lambda \\ \lambda \in \{\alpha, \beta\}}} \left[ \Gamma_{kp}^{\alpha \lambda} n_{\lambda p} (1 - n_{\alpha k}) - \Gamma_{pk}^{\lambda \alpha} n_{\alpha k} (1 - n_{\lambda p}) \right]
% \end{align}

\subsection{Osnovni modeli kopeli}

\subsection{Limita neskončne temperature}
\subsection{Fononske sklopitve}
Tu sta dve (tri) opcije. Najprej pogledamo, kako modificiran (Peierls tip) hopping integral in modificiran orbital energy vplivata na sklopitev. Model spremenimo: atoma znotraj o.c. sta na enakem položaju in se deformirata vedno v isto smer, torej $\vec{u}_{a,i}=\vec{u}_{b,i}=\vec{u}_{i}$:
\begin{align}
  H(\{\mathbf{u}\}) = \underbrace{\sum_{i, \gamma, \tau} \epsilon_{\gamma\tau}(\mathbf{u}_i) \, c_{\gamma,i}^\dagger c_{\tau,i}}_{\text{on-site in CEF}} + \underbrace{\sum_{\langle i,j \rangle, \gamma, \tau} t_{\gamma\tau}(\mathbf{r}_{ij}) \, c_{\gamma,i}^\dagger c_{\tau,j}}_{\text{hopping, med oribitalami}} + \underbrace{V \sum_j n_{a,j} n_{b,j}}_{\text{Coulomb}}
\end{align}
Predpostavimo, da se z deformacijo energija na orbitalo spremeni (podobnokot Kittel, str. 131), torej:
\begin{equation}
\begin{pmatrix}  \epsilon_a^{(0)} & 0 \\ 0 &  \epsilon_b^{(0)} \end{pmatrix} \to \begin{pmatrix}  \epsilon_a^{(0)} + D_a \displaystyle\sum_{\langle ij \rangle} \hat{\mathbf{R}}_{ij}^0 \cdot (\mathbf{u}_j - \mathbf{u}_i) & D_{ab} \displaystyle\sum_{\langle ij \rangle} (\hat{R}_{ij}^0)_x \, \delta \vec{R}_{ij} \\  D_{ab} \displaystyle\sum_{\langle ij \rangle} (\hat{R}_{ij}^0)_x \, \delta \vec{R}_{ij} & \epsilon_b^{(0)} + D_b \displaystyle\sum_{\langle ij \rangle} \hat{\mathbf{R}}_{ij}^0 \cdot (\mathbf{u}_j - \mathbf{u}_i)  \end{pmatrix}
\end{equation}

\begin{align}
\vec r_{\ell,j}=\vec R_j+\vec u_{\ell,j}.
\end{align}

\begin{align}
\vec U_j=\frac{\vec u_{a,j}+\vec u_{b,j}}{2},
\qquad
\vec r_j=\vec u_{a,j}-\vec u_{b,j}.
\end{align}

\begin{align}
\vec u_{a,j}=\vec U_j+\frac{\vec r_j}{2},
\qquad
\vec u_{b,j}=\vec U_j-\frac{\vec r_j}{2}.
\tag{3}
\end{align}
Razvijemo
\begin{align}
H_e[\vec u]
=
H_e[0]+H_{e\text{-}ph}^{(1)}+O(u^2),
\qquad
H_{e\text{-}ph}^{(1)}
=
\left.
\sum_{\ell,j}
\vec u_{\ell,j}\cdot
\frac{\partial H_e}{\partial\vec u_{\ell,j}}
\right|_0.
\end{align}
\begin{align}
  t_{ij}^{(aa)}(\mathbf{r}_{ij}) &\approx t_{ij}^{(a)0} + \alpha_t^{(a)} \, \hat{\mathbf{R}}_{ij}^0 \cdot (\mathbf{u}_j - \mathbf{u}_i)\\
  t_{ij}^{(bb)}(\mathbf{r}_{ij}) &\approx t_{ij}^{(b)0} + \alpha_t^{(b)} \, \hat{\mathbf{R}}_{ij}^0 \cdot (\mathbf{u}_j - \mathbf{u}_i)\\
  t_{ij}^{(ab)}(\mathbf{r}_{ij}) &\approx \underbrace{0}_{\text{ni hoppinga v nedeform. mrezi}} + \, \boldsymbol{\alpha}_t^{(ab)} \cdot (\mathbf{u}_j - \mathbf{u}_i)
\end{align}

\begin{align}
  \boldsymbol{\alpha}_t^{(ab)} = \left. \nabla_{\mathbf{r}} t_{ij}^{(ab)} \right\vert{}_{\mathbf{R}_{ij}^0} \neq 0
\end{align}

\begin{equation}
H_e = \sum_j \left( \epsilon_a c_{ja}^\dagger c_{ja} + \epsilon_b c_{jb}^\dagger c_{jb} \right) - \sum_{\langle i, j \rangle} \left( t_a(\vec{R}_j - \vec{R}_i) c_{ja}^\dagger c_{ia} + t_b(\vec{R}_j - \vec{R}_i) c_{jb}^\dagger c_{ib} + \text{h.c.} \right) + \sum_{\langle i, j \rangle} t_{ab}(\vec{R}_j - \vec{R}_i) c_{ja}^\dagger c_{ib} + \text{h.c.}
\end{equation}

\begin{equation}
\vec{d}_{ij, ab} = \vec{r}_{j,b} - \vec{r}_{i,a} = (\vec{R}_j - \vec{R}_i) + (\vec{r}_{jb} - \vec{r}_{ia}) + (\vec{u}_{j,b} - \vec{u}_{i,a})
\end{equation}

\begin{equation}
t_a = t_a(d_{ij, aa}), \quad t_b = t_b(d_{ij, bb})
\end{equation}

\begin{equation}
t_a = t_a(d_{ij, aa}^0) + \vec{\nabla}_{\vec{r}} t_a(\vec{r}) \Big|_{\vec{d}_{ij, aa}^0} \cdot (\vec{u}_{j,a} - \vec{u}_{i,a})
\end{equation}

\begin{equation}
t_a(\vec{R}_{ij}^0) \approx t_a(R_{ij}^0) + \vec{\nabla}_{\vec{r}} t_a(\vec{r}) \cdot (\vec{u}_j - \vec{u}_i)
\end{equation}

\begin{equation}
t_{ab}(\vec{R}_{ij}^0) \simeq 0 + \vec{\nabla}_{\vec{r}} t_{ab}(\vec{r}) \Big|_{\vec{R}_{ij}^0} \cdot (\vec{u}_j - \vec{u}_i) \approx \alpha_{ab} \hat{e}_{ij} \cdot (\vec{u}_j - \vec{u}_i)
\end{equation}

\begin{equation}
\vec{\nabla} [t_{ab}(x,y)] = \vec{\nabla}(x y \cdot f(r))
\end{equation}

\begin{equation}
H_e = \sum_i \left( \epsilon_a^{(0)} c_{ai}^\dagger c_{ai} + \epsilon_b^{(0)} c_{bi}^\dagger c_{bi} \right) - \sum_{\langle i, j \rangle} \left( t_a c_{aj}^\dagger c_{ia} + t_b c_{bj}^\dagger c_{ib} + \text{h.c.} \right) + V \sum_i n_{ai} n_{bi}
\end{equation}

\begin{equation}
H_{ep} = \sum_{\langle i, j \rangle} \left( \alpha_a c_{ja}^\dagger c_{ia} \hat{R}_{ij}^0 + \alpha_b c_{jb}^\dagger c_{ib} \hat{R}_{ij}^0 + \alpha_{ab} c_{ja}^\dagger c_{ib} + \alpha_{ba} c_{jb}^\dagger c_{ia} + \text{h.c.} \right) \cdot (\vec{u}_j - \vec{u}_i)
\end{equation}

\begin{equation}
V(\vec{R}_i - \vec{R}_j) \simeq V_{\text{onsite}} \delta(\vec{R}_j - \vec{R}_i) + \mathcal{O}(|\vec{u}_j - \vec{u}_i|)
\end{equation}

\begin{equation}
\vec{u}_j = \frac{1}{\sqrt{N_{\text{uc}}}} \sum_{\vec{q}, \lambda} \frac{1}{\sqrt{2M\omega_{\vec{q}\lambda}}} \vec{\xi}_{\vec{q}\lambda} e^{i\vec{q}\cdot\vec{R}_j} \left( b_{\vec{q}\lambda} + b_{-\vec{q}\lambda}^\dagger \right)
\end{equation}

\begin{equation}
c_{aj} = \frac{1}{\sqrt{N_{\text{uc}}}} \sum_{\vec{k}} e^{i\vec{k}\cdot\vec{R}_j} c_{a\vec{k}}
\end{equation}



\begin{equation}
\hat{\delta} = (\delta_x, \delta_y) \in \{\pm \hat{x}, \pm \hat{y}\}
\end{equation}

\begin{equation}
\sum_{\hat{\delta}} \hat{\delta} \cdot (\vec{u}_{i+\delta} - \vec{u}_i) = \hat{x} \cdot (\vec{u}_{i+x} - \vec{u}_i) - \hat{x} \cdot (\vec{u}_{i-x} - \vec{u}_i) + \hat{y} \cdot (\vec{u}_{i+y} - \vec{u}_i) - \hat{y} \cdot (\vec{u}_{i-y} - \vec{u}_i) \approx 2 \vec{\nabla} \cdot \vec{u}
\end{equation}

\begin{equation}
\sum_{\hat{\delta}} \tilde{\delta} \cdot (\vec{u}_{i+\delta} - \vec{u}_i) = (u_{i+x}^y - u_{i-x}^y) - (u_{i+y}^x - u_{i-y}^x)
\end{equation}

\begin{equation}
H_{e-p} = \sum_{i, \hat{\delta}} \left( D_a c_{ai}^\dagger c_{ai} + D_b c_{bi}^\dagger c_{bi} + \alpha_a c_{a, i+\delta}^\dagger c_{ai} + \alpha_b c_{b, i+\delta}^\dagger c_{bi} \right) \hat{\delta} \cdot (\vec{u}_{i+\delta} - \vec{u}_i) + \sum_{i, \hat{\delta}} \left( D_{ab} c_{ai}^\dagger c_{bi} + \text{h.c.} \right) \hat{\delta} \cdot (\vec{u}_{i+\delta} - \vec{u}_i)
\end{equation}

\begin{equation}
\vec{u}_{i+\delta} - \vec{u}_i = \frac{1}{\sqrt{N}} \sum_{\vec{q}, \lambda} \frac{1}{\sqrt{2M\omega_{\vec{q}\lambda}}} \vec{\epsilon}_{\vec{q}\lambda} \left( e^{i\vec{q}\cdot\hat{\delta}} - 1 \right) e^{i\vec{q}\cdot\vec{R}_i^0} \left( b_{\vec{q}\lambda} + b_{-\vec{q}\lambda}^\dagger \right)
\end{equation}


\begin{equation}
H_{e-p} = \frac{1}{\sqrt{N}} \sum_{\vec{k}, \vec{q}, \lambda} \left[ \frac{1}{\sqrt{2M\omega_{\vec{q}\lambda}}} \sum_{\hat{\delta}} \hat{\delta} \cdot \vec{\epsilon}_{\vec{q}\lambda} \left( e^{i\vec{q}\cdot\hat{\delta}} - 1 \right) e^{i\vec{k}\cdot\hat{\delta}} \right] c_{\alpha, \vec{k}+\vec{q}}^\dagger c_{\beta\vec{k}} \left( b_{\vec{q}\lambda} + b_{-\vec{q}\lambda}^\dagger \right)
\end{equation}

\begin{equation}
\sum_{\hat{\delta}} \left( e^{i\vec{k}\cdot\hat{\delta}} - 1 \right) e^{i\vec{k}\cdot\hat{\delta}} \hat{\delta} \cdot \vec{\epsilon}_{\vec{q}\lambda} = \epsilon_{\vec{q}\lambda}^x \left[ \left( e^{iq_x} - 1 \right) - \left( e^{-iq_x} - 1 \right) \right] e^{ik_x}
\end{equation}

\begin{equation}
g_{aa}^\lambda(\vec{k}, \vec{q}) = \frac{2i}{\sqrt{2M}} \frac{1}{\sqrt{\omega_{\vec{q}\lambda}}} \left[ \epsilon_{\vec{q}\lambda}^x \left( D_a \sin(q_x) + \alpha_a \left( \sin(k_x + q_x) - \sin(k_x) \right) \right) + \epsilon_{\vec{q}\lambda}^y (x \to y) \right]
\end{equation}

\begin{equation}
g_{ab}^\lambda(\vec{k}, \vec{q}) = \frac{2i}{\sqrt{2M}} \frac{1}{\sqrt{\omega_{\vec{q}\lambda}}} \left[ \epsilon_{\vec{q}\lambda}^y \left( D_{ab} \sin(q_y) \right) + \epsilon_{\vec{q}\lambda}^x \left( D_{ab} \sin(q_x) \right) \right]
\end{equation}

\begin{equation}
\hat{R}_{ij}^0 = \frac{1}{\sqrt{N}} \sum_{\vec{k}, \vec{q}} \left( c_{a\vec{k}}^\dagger, c_{b\vec{k}}^\dagger \right) \begin{pmatrix} \alpha_{aa} & \alpha_{ab} \\ \alpha_{ba} & \alpha_{bb} \end{pmatrix} \begin{pmatrix} c_{a, \vec{k}+\vec{q}} \\ c_{b, \vec{k}+\vec{q}} \end{pmatrix}
\end{equation}


\begin{equation}
\rho(\vec{r}) = -e \cdot \psi^\dagger(\vec{r}) \psi(\vec{r})
\end{equation}

\begin{equation}
\psi(\vec{r}) = \sum_{i, \alpha} \phi_\alpha(\vec{r} - \vec{R}_i) c_{\alpha i}
\end{equation}

\begin{equation}
\rho(\vec{q}) = -e \sum_{i,j} \sum_{\alpha, \beta} \int \dd{\vec{r}} \, e^{-i\vec{q}\cdot\vec{r}} \phi_\alpha^*(\vec{r} - \vec{R}_i) \phi_\beta(\vec{r} - \vec{R}_j) c_{\alpha i}^\dagger c_{\beta j} = -e \sum_{i, \alpha, \beta} e^{-i\vec{q}\cdot\vec{R}_i} \langle \phi_\alpha | e^{-i\vec{q}\cdot\vec{r}} | \phi_\beta \rangle c_{\alpha i}^\dagger c_{\beta i} \, \delta_{ij}
\end{equation}

\begin{equation}
\rho(\vec{q}) \simeq -e \sum_i e^{-i\vec{q}\cdot\vec{R}_i} \left( n_{ai} + n_{bi} - i \vec{q} \cdot 2 \langle \phi_a(\vec{r}) | \vec{r} | \phi_b(\vec{r}) \rangle \begin{pmatrix} 0 & 1 \\ 1 & 0 \end{pmatrix} \begin{pmatrix} c_{ai} \\ c_{bi} \end{pmatrix} \right)
\end{equation}

\begin{equation}
\sum_i e^{i\vec{q}\cdot\vec{R}_i} c_{ai}^\dagger c_{bi} = \frac{1}{N} \sum_{\vec{k}} c_{a, \vec{k}+\vec{q}}^\dagger c_{b\vec{k}}
\end{equation}
% \begin{align}
% \dot{n}_{\alpha k} = \sum_{p} \bigg( & \Gamma_{\alpha p \to \alpha k} \, n_{\alpha p} (1 - n_{\alpha k}) - \Gamma_{\alpha k \to \alpha p} \, n_{\alpha k} (1 - n_{\alpha p}) + \\
% & + \Gamma_{\beta p \to \alpha k} \, n_{\beta p} (1 - n_{\alpha k}) - \Gamma_{\alpha k \to \beta p} \, n_{\alpha k} (1 - n_{\beta p}) \bigg)
% \end{align}

% \begin{align}
% \Gamma_{\mu p \to \nu k} = \frac{2\pi}{V} \sum_{\lambda} \left| G_{\mu \nu}^{\lambda} (k, q) \right|^2 \left\{ (1 + N_{q}) \delta (E_{\nu k} - E_{\mu p} + \omega_{q \lambda}) + N_{-q} \delta (E_{\nu k} - E_{\mu p} - \omega_{-q \lambda}) \right\}
% \end{align}

% \begin{align}
% \Gamma_{\nu k \to \mu p} = \frac{2\pi}{V} \sum_{\lambda} \left| G_{\mu \nu}^{\lambda} (k, q) \right|^2 \left\{ (1 + N_{-q}) \delta (E_{\mu p} - E_{\nu k} + \omega_{-q \lambda}) + N_{q} \delta (E_{\mu p} - E_{\nu k} - \omega_{q \lambda}) \right\}
% \end{align}

% \begin{align}
% G^\lambda(\vec{k}, \vec{q}) &= \begin{pmatrix} u_{k+q} & -v_{k+q} \\ v_{k+q} & u_{k+q} \end{pmatrix} \begin{pmatrix} g_{a\lambda}(\vec{q}) & 0 \\ 0 & g_{b\lambda}(\vec{q}) \end{pmatrix} \begin{pmatrix} u_k & v_k \\ -v_k & u_k \end{pmatrix} \\[1.5em]
% g_{\substack{a\\b}\lambda}(\vec{q}) &= -i\ \frac{4\pi e^2}{\sqrt{2M\Omega\omega_{\vec{q}\lambda}}} \cdot \frac{1}{|\vec{q}|^2 + k_{\text{TF}}^2} \left( \vec{q} \cdot \vec{\varepsilon}_{\vec{q}\lambda}^* \right)
% \end{align}


\subsection{Fotonske sklopitve v resonatorju}
Začnemo z minimalnim modelom za elektrone: 
\begin{align}
  H_e = H_e^{(0)} + H_V = \sum_k \epsilon_{a,k} c_{a,k}^\dagger c_{a,k} + \sum_k \epsilon_{b,k} c_{b,k}^\dagger c_{b,k} + V \sum_j c_{a,j}^\dagger c_{a,j} c_{b,j}^\dagger c_{b,j}
\end{align}
in privzamemo 2D plast, kjer sta dve orbitali na o.c., ki hibridizirata samo preko Coulombske interakcije. Disperzije za model neinteragirajoč sistem so v tem primeru enake \eqref{free el tb}. Interakcija pa je v prvem redu oblike: 
\begin{equation}
H_{e\gamma}^{(1)} = -\int \dd{\vec{r}} \, \vec{j}_e \cdot \vec{A} = -\int \dd{\vec{r}} \sum_{\vec{q}} {\vec{A}}({\vec{q}}) e^{i\vec{q}\cdot\vec{r}} \vec{j}_e(\vec{r})  = -\sum_{\vec{q}} {\vec{A}}({\vec{q}}) \cdot \vec{j}_e(\vec{-q})\,,
\end{equation}
kjer je 
\begin{align}
  \vec{A}(\vec{q}) = \sum_\lambda \frac{1}{\sqrt{2\mathcal{V}\epsilon_0 \omega_{\vec{q}, \lambda}}} \bm{\epsilon}_{\vec{q}, \lambda} \left(a_{\vec{q}, \lambda}^{\vphantom{\dagger}} + a_{-\vec{q}, \lambda}^\dagger \right) 
\end{align}
V t.b. modelu dobimo tok iz kontinuitetne in Heisenbergove enačbe: 
\begin{align}
  \frac{\partial \rho(\vec{r})}{\partial t} + \vec{\nabla} \cdot \vec{j} = 0 \, \to\,  i\vec{q} \cdot \vec{j}(\vec{q}) = -\frac{\partial \rho(\vec{q})}{\partial t} 
\end{align}
(+ je od IBP) in 
\begin{equation}
\vec{q}\cdot \vec{j}_e(\vec{q}) = i\frac{\partial \rho(\vec{q})}{\partial t} = -[H_e, \rho(\vec{q})]
\end{equation}
Elektronska gostota je
\begin{equation}
\rho(\vec{r}) = -e \cdot \psi^\dagger(\vec{r}) \psi(\vec{r})
\end{equation}
operatorje polja razvijemo po lokaliziranih / Wannierovih orbitalah
\begin{equation}
\psi(\vec{r}) = \sum_{i,\gamma=\{a, b\}} \phi_\gamma(\vec{r} - \vec{R}_i) c_{\gamma i}\,.
\end{equation}
Transformiramo s FT, uvedemo $\vec{r} = \vec{R}_i + \vec{x}$ in razvijemo do dipolov $e^{-i\vec{q}\cdot\vec{x}} \simeq 1 -i\vec{q}\cdot\vec{x}$ (delamo se da se orbitale ne prekrivajo izven o.c.): 
\begin{equation}
\rho(\vec{q}) = -e \sum_{i,j} \sum_{\gamma, \tau} \int \dd{\vec{r}} \, e^{-i\vec{q}\cdot\vec{r}} \phi_\gamma^*(\vec{r} - \vec{R}_i) \phi_\tau(\vec{r} - \vec{R}_j) c_{\gamma i}^\dagger c_{\tau j} = -e \sum_{i,\gamma,\tau} e^{-i\vec{q}\cdot\vec{R}_i} \langle \phi_\gamma | e^{-i\vec{q}\cdot\vec{x}} | \phi_\tau \rangle c_{\gamma i}^\dagger c_{\tau i} 
\end{equation}
Zdaj pa recimo da je dipolni moment samo med različnima orbitalama v o.c. in $n_{\gamma i} = \langle\phi_\gamma |\phi_\gamma\rangle c_{\gamma,i}^\dagger c_{\gamma,i}^ {\vphantom{\dagger}}$:
\begin{equation}
\rho(\vec{q}) \simeq -e \sum_i e^{-i\vec{q}\cdot\vec{R}_i} \left( n_{ai} + n_{bi} - i \vec{q} \cdot \vec{d} \left( c_{a,i}^\dagger c_{b,i}^{\vphantom{\dagger}} + c_{b,i}^\dagger c_{a,i}^{\vphantom{\dagger}} \right) \right)
\end{equation}
zdaj uporabimo zvezo
\begin{equation}
\sum_i e^{i\vec{q}\cdot\vec{R}_i} c_{ai}^\dagger c_{bi} = \sum_{\vec{k}} c_{a, \vec{k}+\vec{q}}^\dagger c_{b\vec{k}}
\end{equation}
z negativnim argumentom je potem elektronska gostota do dipolnega člena (upoštevajoč $\vec{d} = \langle \phi_a(\vec{x}) | \vec{x} | \phi_b(\vec{x}) \rangle$):
\begin{equation}
\rho(-\vec{q}) = -e \sum_{\vec{k}} \begin{pmatrix} c_{a, \vec{k}+\vec{q}}^\dagger & c_{b, \vec{k}+\vec{q}}^\dagger \end{pmatrix} \left( \mathds{1} + i\, \vec{q} \cdot \vec{d} \begin{pmatrix} 0 & 1 \\ 1 & 0 \end{pmatrix} \right) \begin{pmatrix} c_{a\vec{k}} \\ c_{b\vec{k}} \end{pmatrix}
\end{equation}
Zdaj poračunamo komutator, da dobimo tok. \textbf{Lahko vzamemo $H_e^0$, ker za local $V$ se cleni z $\Delta_k \to \Delta$ pokrajsajo!!}, načeloma pa (nisem še preveru, z dipolnim členom ne komutira). Delam z vektorji $\psi_k = \begin{pmatrix} c_{a, \vec{k}} & c_{b, \vec{k}}\end{pmatrix}^\mathsf{T}$: 



% \begin{equation}
% H_e = \sum_i \left( \epsilon_a^{(0)} c_{ai}^\dagger c_{ai} + \epsilon_b^{(0)} c_{bi}^\dagger c_{bi} \right) - \sum_{\langle i, j \rangle} \left( t_a c_{aj}^\dagger c_{ai} + t_b c_{bj}^\dagger c_{bi} + \text{h.c.} \right) + V \sum_i n_{ai} n_{bi}
% \end{equation}

% \begin{equation}
% H_e^0 = \sum_{\vec{k}, \gamma \in \{a,b\}} \varepsilon_{\vec{k}\gamma} c_{\vec{k}\gamma}^\dagger c_{\vec{k}\gamma}
% \end{equation}
\begin{equation}
[H_e^{(0)}, \rho(\vec{q})] =  -e \sum_{\vec{k}, \vec{k'}} \left[ \psi_{\vec{k}}^\dagger \begin{pmatrix} \varepsilon_{a,\vec{k}} & 0 \\ 0 & \varepsilon_{b,\vec{k}} \end{pmatrix} \psi_{\vec{k}}, \psi_{\vec{k'} - \vec{q}}^\dagger \begin{pmatrix} 1 & -i\vec{q}\cdot\vec{d} \\ -i\vec{q}\cdot\vec{d} & 1 \end{pmatrix} \psi_{\vec{k'}} \right]
\end{equation}
uporabimo zvezo:
\begin{align}
  \left[ \psi_{\vec{k}}^\dagger A \psi_{\vec{k}}, \psi_{\vec{k}'-\vec{q}}^\dagger B \psi_{\vec{k}'} \right] &= \psi_{\vec{k}'-\vec{q}}^\dagger \Big( A(\vec{k}'-\vec{q}) B(\vec{k}') - B(\vec{k}'-\vec{q}) A(\vec{k}') \Big) \psi_{\vec{k}'} , \delta_{\vec{k},  \vec{k}'-\vec{q}} \\ &= \psi_{\vec{k}}^\dagger \Big( A(\vec{k}) B(\vec{k}+\vec{q}) - B(\vec{k}) A(\vec{k}+\vec{q}) \Big) \psi_{\vec{k}+\vec{q}}
\end{align}
od tod 
\begin{equation}
\vec{q} \cdot \vec{j}_e(\vec{q}) = -[H_e^{(0)}, \rho(\vec{q})] = -e \sum_{\vec{k}} \psi_{\vec{k}}^\dagger  \begin{pmatrix} \varepsilon_{a,\vec{k}+\vec{q}} - \varepsilon_{a,\vec{k}} & i(\vec{q}\cdot\vec{d})\left(\varepsilon_{a,\vec{k}} - \varepsilon_{b,\vec{k}+\vec{q}}\right) \\ i(\vec{q}\cdot\vec{d})\left(\varepsilon_{b,\vec{k}} - \varepsilon_{a,\vec{k}+\vec{q}}\right) & \varepsilon_{b,\vec{k}+\vec{q}} - \varepsilon_{b,\vec{k}} \end{pmatrix}  \psi_{\vec{k}+\vec{q}}
\end{equation}
da izrazimo tok, naredimo približek, da prispevajo le majhni $\vec{q}\to 0$, zato:
\begin{equation}
\epsilon_{a\vec{k}} - \epsilon_{a, \vec{k}+\vec{q}} \simeq - \vec{q} \cdot \vec{\nabla}_{\vec{k}} \epsilon_{\vec{k}+\vec{q}/2}\,,
\end{equation}
kar drži točno za proste elektrone in približno za tight binding disperzije. 
% Note (baje): 
% For any tight-binding dispersion $\varepsilon_{\mathbf{k}}$, the difference $\varepsilon_{\mathbf{k}} - \varepsilon_{\mathbf{k}-\mathbf{q}}$ can always be written exactly (without taking derivatives) as an inner product:$$\varepsilon_{\mathbf{k}} - \varepsilon_{\mathbf{k}-\mathbf{q}} = \mathbf{q} \cdot \boldsymbol{\mathcal{V}}(\mathbf{k}, \mathbf{q})$$where $\boldsymbol{\mathcal{V}}(\mathbf{k}, \mathbf{q})$ is the finite-momentum velocity kernel defined as:$$\boldsymbol{\mathcal{V}}(\mathbf{k}, \mathbf{q}) = \int_0^1 d\lambda \, \boldsymbol{\nabla}_{\mathbf{k}} \varepsilon_{\mathbf{k} - (1-\lambda)\mathbf{q}}$$ ... ce ne expandamo za $\vec{q}\to 0$. 
In general ce ne expandamo: 
\begin{align}
  \mathbf{j}_e(\mathbf{q}) = -\frac{e\mathbf{q}}{\Vert{}\mathbf{q}\Vert{}^2} \sum_{\mathbf{k}} \psi_\mathbf{k}^\dagger {M}(\mathbf{k}, \mathbf{q}) \psi_{\mathbf{k}+\mathbf{q}} \;+\; \mathbf{q} \times \mathbf{C}(\mathbf{q})
\end{align}
kjer so $\mathbf{C}(\mathbf{q})$ arbitrarna vektorska polja (tokovi brez divergence, magnetizacijski tokovi). Ce pa naredimo ta priblizek pa dobimo za tok v splošnem, indekse zamaknemo za $-\vec{q}$: 
% \begin{equation}
% i\vec{q} \cdot \vec{j}(\vec{q}) \simeq i \cdot \left( -e \sum_{\vec{k}} \psi_{\vec{k}}^\dagger \begin{pmatrix} -\vec{q}\cdot\vec{\nabla}_{\vec{k}} \epsilon_{a, \vec{k}+\vec{q}/2} & -i(\vec{q}\cdot\vec{d})(\epsilon_{b, \vec{k}+\vec{q}} - \epsilon_{a\vec{k}}) \\ i(\vec{q}\cdot\vec{d})(\epsilon_{b, \vec{k}+\vec{q}} - \epsilon_{a\vec{k}}) & -\vec{q}\cdot\vec{\nabla}_{\vec{k}} \epsilon_{b, \vec{k}+\vec{q}/2} \end{pmatrix} \psi_{\vec{k}+\vec{q}} \right)
% \end{equation}
\begin{tcolorbox}
\begin{equation}
\vec{j}_e(\vec{q}) \simeq -e \sum_{\vec{k}} \psi_{\vec{k}-\vec{q}}^\dagger \begin{pmatrix} \vec{\nabla}_{\vec{k}} \epsilon_{a,\vec{k}-\vec{q}/2} & i\vec{d}(\epsilon_{a, \vec{k}-\vec{q}} - \epsilon_{b\vec{k}}) \\ i\vec{d}(\epsilon_{b, \vec{k}-\vec{q}} - \epsilon_{a\vec{k}}) & \vec{\nabla}_{\vec{k}} \epsilon_{b,\vec{k}-\vec{q}/2} \end{pmatrix} \psi_{\vec{k}}\,.
\end{equation}
\end{tcolorbox}
Od tod splošna Hamiltonka v prvem redu:
\begin{tcolorbox}
\begin{align}
H_{e\gamma}^{(1)}
&= -\sum_{\vec{q}} \vec{j}_e(\vec{-q}) \cdot \vec{A}(\vec{q}) = \\
&= \sum_{\vec{k},\vec{q},\lambda} \frac{e}{\sqrt{2\mathcal{V}\varepsilon_0\omega_{\vec{q}, \lambda}}}\, \psi_{\vec{k}+\vec{q}}^\dagger \underbrace{ \begin{pmatrix} \bm{\epsilon}_{\vec{q}\lambda}\cdot \vec{v}_a\!\left(\vec{k}+\frac{\vec{q}}{2}\right) & i\bigl(\bm{\epsilon}_{\vec{q}\lambda}\cdot\vec{d}\bigr) \bigl(\varepsilon_{a,\vec{k}+\vec{q}}-\varepsilon_{b,\vec{k}}\bigr) \\ i\bigl(\bm{\epsilon}_{\vec{q}\lambda}\cdot\vec{d}\bigr) \bigl(\varepsilon_{b,\vec{k}+\vec{q}}-\varepsilon_{a,\vec{k}}\bigr) & \bm{\epsilon}_{\vec{q}\lambda}\cdot \vec{v}_b\!\left(\vec{k}+\frac{\vec{q}}{2}\right) \end{pmatrix} }_{\displaystyle M^\lambda(\vec{k},\vec{q})} \psi_{\vec{k}} \left(a_{\vec{q}\lambda}+a_{-\vec{q}\lambda}^\dagger\right).
\end{align}
\end{tcolorbox}

Za setup od Curtisa: material 2D v $xy$ ravnini, stoječe EM valovanje v $z$, vzamemo osnovni mode ($n=1$), ki prispeva največ, v tem primeru sta dve možni polarizaciji:
\begin{equation}
\vec{\bm{\epsilon}}_{\text{TE}}(z) = \left( \hat{\mathbf{z}} \times \hat{\mathbf{q}} \right) \sin\left(\frac{\pi}{L} z\right)
\end{equation}
\begin{equation}
\vec{\bm{\epsilon}}_{\text{TM}}(z) = \frac{\omega_0}{\omega_{\vec{q}}} \hat{\mathbf{q}} \sin\left(\frac{\pi}{L} z\right) + i \frac{c\, q}{\omega_{\vec{q}}} \hat{\mathbf{z}} \cos\left(\frac{\pi}{L} z\right)
\end{equation}
kjer je $\vec{q} = q(\cos\varphi, \sin\varphi)$ valovni vektor v ravnini, $\omega_0 = \pi c / L$ in $\omega_{\vec{q}} = \sqrt{\omega_0^2 + c^2 q^2}$.
Na sredini resonatorja, pri $z = L/2$, je $\sin(\pi z / L) = 1$ in $\cos(\pi z / L) = 0$, zato:
\begin{equation}
\vec{\bm{\epsilon}}_{\text{TE}}\left(\tfrac{L}{2}\right) = \hat{\mathbf{z}} \times \hat{\mathbf{q}} = \begin{pmatrix} -\sin\varphi \\ \cos\varphi \\ 0 \end{pmatrix}
\end{equation}
\begin{equation}
\vec{\bm{\epsilon}}_{\text{TM}}\left(\tfrac{L}{2}\right) = \frac{\omega_0}{\omega_{\vec{q}}} \hat{\mathbf{q}} = \frac{\omega_0}{\omega_{\vec{q}}} \begin{pmatrix} \cos\varphi \\ \sin\varphi \\ 0 \end{pmatrix}
\end{equation}
Iz tega imamo pri $z = L/2$ možne sklopitve glede na model ($\vec{d} = d(\cos\theta, \sin\theta)$ v ravnini oz. $\vec{d} = d\,\hat{\mathbf{z}}$ izven ravnine):

\begin{table}[h!]
\centering
\begin{tabular}{|l|l|l|l|}
  \hline
\textbf{Mode} & $\vec{\bm{\epsilon}} \cdot \vec{v}_\gamma\left(\vec{k} + \frac{\vec{q}}{2}\right)$ & $\vec{\bm{\epsilon}} \cdot \vec{d}$, $\vec{d}$ in-plane & $\vec{\bm{\epsilon}} \cdot \vec{d}$, $\vec{d} \parallel \hat{\mathbf{z}}$ \\
  \hline

\textbf{TE} & $2 t_\gamma \left( -\sin\varphi \sin\left(k_x + \frac{q_x}{2}\right) + \cos\varphi \sin\left(k_y + \frac{q_y}{2}\right) \right)$ & $d \sin(\theta - \varphi)$ & $0$ \\
  \hline

\textbf{TM} & $\frac{\omega_0}{\omega_{\vec{q}}} 2 t_\gamma \left( \cos\varphi \sin\left(k_x + \frac{q_x}{2}\right) + \sin\varphi \sin\left(k_y + \frac{q_y}{2}\right) \right)$ & $\frac{\omega_0}{\omega_{\vec{q}}} d \cos(\theta - \varphi)$ & $i\, d\, \frac{c\, q}{\omega_{\vec{q}}} \cos\left(\frac{\pi}{L} z\right)$ ($= 0$ pri $z = \frac{L}{2}$) \\
  \hline

\end{tabular}
\end{table}
dipolni moment je neničeln za liho-sode kombinacije orbital in plane, tega gledamo samo glede na o.c.:
\begin{equation}
\vec{d} = \langle \phi_a | \vec{x} | \phi_b \rangle = \int \phi_a^*(\vec{x}) \, \vec{x} \, \phi_b(\vec{x}) \, \dd{\vec{x}}
\end{equation}

Kaj pa če gledamo namesto tight binding, proste elektrone, torej 
\begin{equation}
\epsilon_{\gamma \vec{k}} = \pm \frac{G}{2} \mp 2t (\cos k_x + \cos k_y) \quad \text{za}\quad k_x, k_y \sim 0
\end{equation}
od tod
\begin{equation}
\epsilon_{\gamma \vec{k}} \approx \pm \frac{G}{2} \mp 2t \left( 1 - \frac{k_x^2}{2} + 1 - \frac{k_y^2}{2} \right) = \pm \left( \frac{G}{2} - 4t \right) \pm t (k_x^2 + k_y^2) = \pm (\epsilon^0 + t \vec{k} \cdot \vec{k})
\end{equation}
Torej je Hamiltonka
\begin{equation}
H_e = \sum_{\gamma, \vec{k}} \left( \epsilon_{\gamma}^0+ t_\gamma \vec{k}^2 \right) c_{\gamma,\vec{k}}^\dagger c_{\gamma,\vec{k}}
\end{equation}
Na drug način lahko poračunamo tok, ker razvijemo
\begin{equation}
\psi(\vec{r}) = \frac{1}{\sqrt{\mathcal{V}}}\sum_{\vec{k},\gamma=\{a, b\}}  e^{i \vec{k}\cdot \vec{r}} c_{\gamma, \vec{k}}\,.
\end{equation} 
in za tok uporabimo 
\begin{equation}
\vec{j}_e(\vec{r}) = -\frac{e}{2mi} \left( \psi^\dagger(\vec{r}) \vec{\nabla} \psi(\vec{r}) - \vec{\nabla}\psi^\dagger(\vec{r}) \psi(\vec{r}) \right) = -\frac{e}{2mi} \frac{1}{\mathcal{V}} \sum_{\gamma, \vec{k}, \vec{k}'} \left( c_{\gamma \vec{k}'}^\dagger c_{\gamma \vec{k}} e^{-i\vec{k}'\cdot\vec{r}} (i\vec{k}) e^{i\vec{k}\cdot\vec{r}} - c_{\gamma \vec{k}'}^\dagger c_{\gamma \vec{k}} (-i\vec{k}') e^{-i\vec{k}'\cdot\vec{r}} e^{i\vec{k}\cdot\vec{r}} \right)
\end{equation}
s FT dobimo
\begin{equation}
\vec{j}_e(\vec{-q}) = \int \dd{\vec{r}} \, \vec{j}_e(\vec{r}) e^{i\vec{q}\cdot\vec{r}} = -\frac{e}{2m} \sum_{\gamma, \vec{k}} \left( c_{\gamma, \vec{k}+\vec{q}}^\dagger c_{\gamma \vec{k}} \, \vec{k} + (\vec{k}+\vec{q}) c_{\gamma, \vec{k}+\vec{q}}^\dagger c_{\gamma \vec{k}} \right) = -\frac{e}{m} \sum_{\vec{k}} \left( \vec{k} + \frac{\vec{q}}{2} \right) \psi_{\vec{k}+\vec{q}}^\dagger \psi_{\vec{k}}
\end{equation}
kar bi dobili tudi iz bolj splošne relacije. V tem primeru dobimo enako kot Curtis, $m = 1/2t$
\begin{tcolorbox}
\begin{equation}
H_{e\gamma}^{(1)} = \frac{e}{m} \sum_{\vec{k}, \vec{q}, \lambda}  \frac{1}{\sqrt{2\mathcal{V} \epsilon_0 \omega_{\vec{q}, \lambda}}} \left( \vec{k} + \frac{\vec{q}}{2} \right) \cdot \bm{\epsilon}_{\vec{q}\lambda} \, \psi_{\vec{k}+\vec{q}}^\dagger \psi_{\vec{k}} \left( a_{\vec{q}\lambda} + a_{-\vec{q}\lambda}^\dagger \right)
\end{equation}
\end{tcolorbox}

\nocite{*}
\printbibliography
\end{document} 